\section*{الفصل \ref{ch:g8:fronts-and-weather-maps} --- الكتل الهوائية والجبهات وخرائط الطقس}

\begin{solution}{exo:g8:fronts-and-weather-maps:1}
(أ) بحرية مدارية: دافئة رطبة. (ب) قارية قطبية: باردة
جافة. (ج) قارية مدارية: حارة جافة. (د) بحرية قطبية: باردة
رطبة.
\end{solution}

\begin{solution}{exo:g8:fronts-and-weather-maps:2}
خط أزرق عليه مثلثات: \omterm{def:g8:fronts-and-weather-maps:front}{جبهة باردة}. خط أحمر عليه أنصاف دوائر: \omterm{def:g8:fronts-and-weather-maps:front}{جبهة
دافئة}. خط يصل بين أماكن متساوية الضغط: \omterm{def:g8:fronts-and-weather-maps:isobar}{خط تساوي الضغط}. ع: مركز
\omterm{def:g8:fronts-and-weather-maps:depression}{مرتفع جوي} (ضغط عالٍ). م: مركز \omterm{def:g8:fronts-and-weather-maps:depression}{منخفض جوي}
(ضغط منخفض).
\end{solution}

\begin{solution}{exo:g8:fronts-and-weather-maps:3}
(أ) صواب. (ب) خطأ: فهي تصل قبلها، بعيدًا جدًا أمامها، عاليًا في السماء.
(ج) خطأ: \omterm{def:g8:fronts-and-weather-maps:isobar}{خطوط تساوي الضغط} المتقاربة تعني ريحًا قوية. (د) صواب.
\end{solution}

\begin{solution}{exo:g8:fronts-and-weather-maps:4}
(أ) \qty{980}{hPa}. (ب) \qty{1034}{hPa}. (ج) \qty{4}{hPa}. (د) $1034 -
980 = \qty{54}{hPa}$.
\end{solution}

\begin{solution}{exo:g8:fronts-and-weather-maps:5}
خطوط رقيقة من \omterm{def:g4:clouds-rain-snow:cloud-families}{السمحاق}؛ حجاب رمادي من \omterm{def:g4:clouds-rain-snow:cloud-families}{الطبقي} المتوسط؛ \omterm{def:g4:clouds-rain-snow:cloud}{السحب} تنخفض
وتزداد سماكة؛ \omterm{def:g4:clouds-rain-snow:precipitation}{مطر} متواصل؛ هواء أدفأ وأرطب.
\end{solution}

\begin{solution}{exo:g8:fronts-and-weather-maps:6}
عند \omterm{def:g8:fronts-and-weather-maps:front}{الجبهة الباردة}، يندفع الهواء البارد الكثيف تحت الهواء الدافئ كأنه
\emph{إسفين}: فيرتفع الهواء الدافئ على منحدر \emph{شديد الانحدار}، ويصعد بسرعة
ويبني سحبًا عالية. وعند \omterm{def:g8:fronts-and-weather-maps:front}{الجبهة الدافئة}، يتسلق الهواء الدافئ الهواء البارد
ببطء على \emph{منحدر لطيف}: فيصعد شيئًا فشيئًا على مسافة
طويلة، وتنتشر سحبه في طبقات.
\end{solution}

\begin{solution}{exo:g8:fronts-and-weather-maps:7}
(أ) قرب م: \omterm{def:g8:fronts-and-weather-maps:isobar}{فخطوط تساوي الضغط} أكثر تقاربًا هناك، والأسهم
أطول. (ب) من الشرق (فالأسهم تشير غربًا، مائلة قليلًا نحو
الجنوب الغربي). (ج) من الشرق (فالأسهم تشير غربًا). (د)
عكس اتجاه عقارب الساعة.
\end{solution}

\begin{solution}{exo:g8:fronts-and-weather-maps:8}
(أ) $1016 - 1003 = \qty{13}{hPa}$. (ب) \omterm{def:g8:fronts-and-weather-maps:depression}{منخفض جوي}: فالضغط يواصل
الانخفاض. (ج) \omterm{def:g4:clouds-rain-snow:cloud}{سحب} وريح \omterm{def:g4:clouds-rain-snow:precipitation}{ومطر}، وعلى الأرجح مع \omterm{def:g8:fronts-and-weather-maps:front}{جبهة}.
\end{solution}

\begin{solution}{exo:g8:fronts-and-weather-maps:9}
(أ) cT، قارية مدارية. (ب) حارة جافة. (ج) فوق \omterm{def:g2:water-on-earth:ocean}{البحر} تحمل بخار
الماء: فتصير أرطب، ويبرّد \omterm{def:g2:water-on-earth:ocean}{البحر} قليلًا الهواء القريب من
الماء.
\end{solution}

\begin{solution}{exo:g8:fronts-and-weather-maps:10}
(أ) $4 - 1 = 3$ أيام. (ب) $3 \div 3 = 1$ يوم لكل عشر سنوات. (ج) يلزم
$5 - 1 = 4$ أيام إضافية: أي نحو 40 سنة.
\end{solution}

\begin{solution}{exo:g8:fronts-and-weather-maps:11}
(أ) نحو \qty{1350}{km} (أربعة أطوال ونصف من \omterm{def:g2:drawing-to-map:scale}{مقياس الرسم}). (ب) \qty{4}{hPa}
على \qty{80}{km}: $4 \div 80 \times 100 = \qty{5}{hPa}$ لكل
\qty{100}{km}. (ج) $4 \div 140 \times 100 \approx \qty{2.9}{hPa}$ لكل
\qty{100}{km}. (د) قرب م، حيث يتغير الضغط أسرع؛ نعم، فالأسهم
أطول قرب م.
\end{solution}

\begin{solution}{exo:g8:fronts-and-weather-maps:12}
(أ) \omterm{def:g8:fronts-and-weather-maps:front}{الجبهة الباردة}: $300 \div 50 = 6$ ساعات. \omterm{def:g8:fronts-and-weather-maps:front}{والجبهة الدافئة} تتحرك بسرعة
\qty{25}{km} في الساعة: $200 \div 25 = 8$ ساعات. \omterm{def:g8:fronts-and-weather-maps:front}{فالجبهة الباردة} تصل
قبل الأخرى بساعتين. (ب) البلدة الأولى: \omterm{def:g4:clouds-rain-snow:cloud}{سحب} شاهقة، وزخات غزيرة، وهبّات،
ثم هواء أبرد وسماء تصفو. والبلدة الثانية: \omterm{def:g4:clouds-rain-snow:cloud}{سحب} طبقية تنخفض
\omterm{def:g4:clouds-rain-snow:precipitation}{ومطر} طويل متواصل، ثم هواء أدفأ وأرطب. (ج) البلدة الثانية،
التي يبلغها الهواء الدافئ.
\end{solution}

\begin{solution}{pb:g8:fronts-and-weather-maps:1}
\textbf{1.} في الهواء الدافئ، بين الجبهتين. \omterm{def:g8:fronts-and-weather-maps:front}{والجبهة الدافئة} قد
مرت بها من قبل.

\textbf{2.} من الجنوب الغربي.

\textbf{3.} نحو الشرق: فمثلثات \omterm{def:g8:fronts-and-weather-maps:front}{الجبهة الباردة} تقع على الجانب
الذي تتحرك نحوه، وهي تشير إلى الشرق، نحو البلدة.

\textbf{4.} \omterm{def:g4:clouds-rain-snow:cloud}{سحب} شاهقة، وزخات غزيرة، وربما رعد \omterm{def:g4:clouds-rain-snow:precipitation}{وبرد}،
وهبّات ريح؛ ثم هواء أبرد، وريح من الغرب أو الشمال الغربي،
وسماء تصفو.

\textbf{5.} على بعد \qty{180}{km} في الساعة \babelsublr{06:00} وعلى بعد \qty{300}{km} في الساعة \babelsublr{09:00}.

\textbf{6.} $300 - 180 = \qty{120}{km}$.

\textbf{7.} $120 \div 3 = \qty{40}{km}$ في الساعة.

\textbf{8.} البلدة على بعد \qty{480}{km}: $480 - 300 = \qty{180}{km}$.

\textbf{9.} $180 \div 40 = 4.5$ ساعة، أي 4 ساعات و30 دقيقة.

\textbf{10.} \babelsublr{09:00} + 4~ساعات و30~دقيقة = \textbf{\babelsublr{13:30}}.

\textbf{11.} لا: فالزخات ستنتهي نحو الساعة \babelsublr{14:30}، قبل
بداية المباراة.

\textbf{12.} $180 \div 30 = 6$ ساعات: أي في الساعة \babelsublr{15:00}، لحظة بداية المباراة تمامًا.
التنبؤ: «ينبغي أن تبلغ \omterm{def:g8:fronts-and-weather-maps:front}{الجبهة الباردة} البلدة نحو الساعة \textbf{\babelsublr{13:30}}،
مع ساعة من الزخات الغزيرة والهبّات، ثم تصفو السماء: والأرجح أن المباراة
ستُلعب. وإذا تباطأت \omterm{def:g8:fronts-and-weather-maps:front}{الجبهة}، فقد تصل في الساعة \babelsublr{15:00} وتؤخر
المباراة.»
\end{solution}
